\section{Independent singularity patterns in random alternating spaces}
\label{sec:alternating}

This section proves the second main ingredient without referring to the
observation model. Its content is joint independence across linearly related
matrices; the one-matrix rank law is classical \cite{fulman,pott}.

\begin{theorem}[Projective independence of nonsingularity]
\label{thm:pencil}
Fix $q$, $k$, and pairwise distinct projective points
$[\alpha_1],\ldots,[\alpha_\ell]\in\mathbb P(\F^k)$. Let
$A_1,\ldots,A_k$ be independent uniform alternating $n\times n$ matrices,
where $n$ tends to infinity through even integers. Set
\[
 B_i=\sum_{j=1}^k\alpha_{ij}A_j,\qquad
 I_i=\ind_{\{B_i\text{ is nonsingular}\}}.
\]
Then $(I_1,\ldots,I_\ell)$ converges in total variation to $\ell$ independent
Bernoulli variables with parameter $\omega_q$.
\end{theorem}

\begin{corollary}[The complete pattern and the binomial law]
\label{cor:random-space}
Let $S_n$ be a uniform $k$-dimensional subspace of $\Alt(\F^n)$, with
$q,k$ fixed. The number of nondegenerate projective forms in $S_n$ converges
in total variation to $\Bin(m_k,\omega_q)$. In coordinates given by a uniform
ordered basis, every fixed subset $T\subseteq\mathbb P(\F^k)$ satisfies
\begin{equation}
 \Pbb\{\text{the nonsingular projective combinations are exactly }T\}
 \longrightarrow \omega_q^{|T|}(1-\omega_q)^{m_k-|T|}.
 \label{eq:pattern}
\end{equation}
In particular every such pattern occurs for all sufficiently large even
$n$. The proportion of $k$-dimensional alternating matrix spaces in which
every nonzero matrix is nonsingular tends to $\omega_q^{m_k}$.
\end{corollary}

\subsection{A tensor-rank rigidity lemma}

\begin{lemma}[Equality forces a direct sum]
\label{lem:tensor}
Let $U_1,\ldots,U_\ell$ be subspaces of a finite-dimensional vector space
$V$ over $\F$, and let $\alpha_i\in\F^k$ be nonzero and pairwise
nonproportional. Define
\[
 U=\sum_iU_i,\qquad W=\sum_i\alpha_i\otimes U_i\subseteq\F^k\otimes V.
\]
Then $\dim W\ge\dim U$, with equality if and only if
$U=\bigoplus_iU_i$.
\end{lemma}

\begin{proof}
Extend scalars to an infinite field containing $\F$. Tensor ranks,
dimensions, and the direct-sum property are unchanged. Choose a linear
functional $f$ on the coefficient space with $f(\alpha_i)\ne0$ for every
$i$. Contraction $C_f=f\otimes\mathrm{id}$ maps $W$ onto $U$, so
$\dim W\ge\dim U$.

If equality holds, this contraction is an isomorphism. Because the
$\alpha_i$ are pairwise nonproportional, a second functional $g$ can be
chosen such that the ratios $g(\alpha_i)/f(\alpha_i)$ are all distinct:
each forbidden equality is a proper linear hyperplane in the space of
functionals, and finitely many proper hyperplanes do not cover an infinite
vector space. The endomorphism $C_gC_f^{-1}$ of $U$ acts on $U_i$ by the
scalar $g(\alpha_i)/f(\alpha_i)$. Eigenspaces for distinct eigenvalues
have a direct sum, proving the forward implication. Conversely, if the
$U_i$ have a direct sum, then the subspaces $\alpha_i\otimes U_i$ also
have a direct sum of the same total dimension.
\end{proof}

\subsection{Mixed kernel-subspace moments}

Let $\nu_i=\dim\ker B_i$. For fixed nonnegative integers
$r_1,\ldots,r_\ell$, put $r=\sum_i r_i$.

\begin{lemma}[Factorization of all fixed mixed moments]
\label{lem:mixed}
As $n\to\infty$ through even integers,
\begin{equation}
 \E\prod_{i=1}^\ell\qbinom{\nu_i}{r_i}
 \longrightarrow
 \prod_{i=1}^\ell\frac{q^{-r_i(r_i-1)/2}}{g_{r_i}}.
 \label{eq:mixed-moments}
\end{equation}
The contribution from tuples of kernel subspaces not forming a direct sum
is $O_{q,k,\ell,\boldsymbol r}(q^{-n})$.
\end{lemma}

\begin{proof}
Expand the product as a sum over subspaces $U_i\le V$ with
$\dim U_i=r_i$, asking for the probability that $U_i\le\ker B_i$
for all $i$.

First suppose the $U_i$ form a direct sum $U$ of dimension $r$. Choose a
basis adapted to these subspaces and a complement. Because the matrices
$A_j$ are uniform alternating matrices, a common change of basis preserves
their independent uniform distribution. Each matrix entry associated to an
unordered pair of distinct basis vectors is a uniform independent vector
of $\F^k$ across the $k$ matrices.
For a pair with one endpoint in $U_i$ and the other in the complement there
is one equation, with coefficient vector $\alpha_i$. For a pair inside
$U_i$ there is one equation. For a pair with endpoints in distinct $U_i,U_j$
there are two independent equations, since $\alpha_i,\alpha_j$ are linearly
independent. Thus the total constraint codimension is exactly
\begin{align}
 r(n-r)+\sum_i\binom{r_i}{2}+2\sum_{i<j}r_ir_j
 &=\sum_i\left(nr_i-\frac{r_i(r_i+1)}2\right).
 \label{eq:direct-codim}
\end{align}
The probability is $q$ to the negative of this number. A tuple of
independent uniform $r_i$-subspaces forms a direct sum with probability
$1-O_{\boldsymbol r,q}(q^{-n})$. This follows, for example, by a union bound
over nonzero vectors in the intersections with the preceding span; that
span has bounded dimension. Therefore the number of direct-sum tuples is
\[
 (1+O(q^{-n}))\prod_i\qbinom n{r_i}.
\]
Multiplying by the constraint probability and using
$\qbinom n{s}=q^{s(n-s)}(1+o(1))/g_s$ gives the right side of
\eqref{eq:mixed-moments}.

For the remaining tuples write $s=\dim\sum_iU_i<r$ and
$w=\dim\sum_i\alpha_i\otimes U_i$. By \cref{lem:tensor}, $w\ge s+1$.
There are $O_{q,\boldsymbol r}(q^{s(n-s)})$ choices for the tuple: choose
its $s$-dimensional span and then a bounded number of subspace configurations
inside that span. The equations obtained by pairing against a fixed basis
vector outside the span have rank $w$. Different outside basis vectors use
disjoint random entries, so these equations alone impose codimension
$(n-s)w$. Hence the total contribution of such tuples is at most a constant
times
\[
 q^{s(n-s)}q^{-(n-s)w}
 \le q^{-(n-s)}=O(q^{-n}).
\]
There are only finitely many possible values of $s$ for the fixed
$\boldsymbol r$, completing the proof.
\end{proof}

The expressions $\qbinom{\nu}{r}$ are polynomials of degree $r$ in
$x=q^\nu$, with nonzero leading coefficient:
\[
 \qbinom{\nu}{r}=
 \prod_{j=0}^{r-1}\frac{x-q^j}{q^r-q^j}.
\]
They form a basis for polynomials in $x$. Consequently
\cref{lem:mixed} gives factorization in the limit for any fixed product of
polynomials in the $q^{\nu_i}$. That observation alone would not justify
convergence of the nonsingularity indicators. The next step supplies the
needed approximation and uses the even parity of every $\nu_i$.

\subsection{Even-nullity interpolation}

\begin{lemma}[One-matrix law and a uniform tail bound]
\label{lem:alternating-marginal}
Let $A$ be uniform alternating of even order $n$. If $s$ is even and
$0\le s\le n$, then
\begin{equation}
 \Pbb(\dim\ker A=s)
 =q^{-s(s-1)/2}\frac{g_n}{g_s e_{(n-s)/2}},\qquad
 e_j=\prod_{i=1}^j(1-q^{-2i}).
 \label{eq:alt-mass}
\end{equation}
It is zero for odd $s$. In particular, for a constant depending only on $q$,
\[
 \Pbb(\dim\ker A=s)\le C_q q^{-s(s-1)/2},
\]
and $\Pbb(A\text{ nonsingular})\to\omega_q$.
\end{lemma}

\begin{proof}
Choose the radical, an $s$-subspace, and choose a nondegenerate alternating
form on the quotient. There are
\[
 \qbinom ns\,
 q^{(n-s)(n-s-1)/2}
 \prod_{j=1}^{(n-s)/2}(1-q^{-(2j-1)})
\]
forms with that radical dimension. Dividing by $q^{n(n-1)/2}$ and cancelling
products proves \eqref{eq:alt-mass}. The denominator is bounded below by
$g_\infty e_\infty>0$, giving the tail bound. At $s=0$ the product is
the finite odd product converging to $\omega_q$.
\end{proof}

For $R\ge1$, define the polynomial
\begin{equation}
 P_R(x)=\prod_{j=1}^R\frac{1-xq^{-2j}}{1-q^{-2j}}.
 \label{eq:interpolation}
\end{equation}
It equals one at $x=1$ and zero at $x=q^2,q^4,\ldots,q^{2R}$.
For even $s\ge2R+2$,
\[
 |P_R(q^s)|\le e_\infty^{-1}q^{Rs-R(R+1)}.
\]
Using \cref{lem:alternating-marginal} and writing $s=2R+2+2h$ gives
\begin{align}
 \sup_{n\text{ even}}
 \E\bigl|P_R(q^{\dim\ker A})-\ind_{\{\ker A=0\}}\bigr|^2
 &\le C_q\sum_{h\ge0}q^{-R-1-3h-2h^2}
 \le C'_q q^{-R}.
 \label{eq:L2-approx}
\end{align}
In particular, the approximation holds uniformly in the matrix order. This
parity-aware bound avoids relying on a generally invalid moment-determinacy
shortcut.

\begin{proof}[Proof of \cref{thm:pencil}]
It suffices to show that for every subset of the $\ell$ indices the
probability that all the corresponding matrices are nonsingular tends to
$\omega_q$ to the size of that subset. Relabel and consider the first $u$
indices. Put $X_i=q^{\nu_i}$ and $I_i=\ind_{\{\nu_i=0\}}$.
For a fixed $R$, expand the difference of products by telescoping and apply
Cauchy--Schwarz to each term:
\begin{align}
 \left|\E\prod_{i=1}^u I_i-\E\prod_{i=1}^uP_R(X_i)\right|
 \le\sum_{i=1}^u
 \bigl(\E|I_i-P_R(X_i)|^2\bigr)^{1/2}
 \left(\E\prod_{j<i}P_R(X_j)^2\right)^{1/2}.
 \label{eq:telescope}
\end{align}
The remaining indicator factors have absolute value at most one and may be
dropped in this bound. By \cref{lem:mixed}, for fixed $R$ the mixed
polynomial expectations factor in the limit $n\to\infty$. The limiting
one-matrix $L^2$ norm of $P_R$ is bounded by $1+O_q(q^{-R/2})$, using
\eqref{eq:L2-approx}. The first factor in each summand is
$O_q(q^{-R/2})$. Thus the limsup of \eqref{eq:telescope} is
$O_{q,u}(q^{-R/2})$.

Also by mixed polynomial factorization,
\[
 \lim_{n\to\infty}\E\prod_{i=1}^uP_R(X_i)
 =\left(\lim_{n\to\infty}\E P_R(q^{\dim\ker A})\right)^u.
\]
The inner limit differs from $\omega_q$ by $O_q(q^{-R/2})$.
Letting $R\to\infty$ proves the desired probability limit.
Inclusion-exclusion recovers each zero-one pattern from these all-one
probabilities. There are finitely many patterns, so convergence is in total
variation.
\end{proof}

\begin{proof}[Proof of \cref{cor:random-space}]
The probability that $k$ independent uniform vectors in the
$D$-dimensional space $E$ are linearly dependent is
\[
 1-\prod_{j=0}^{k-1}(1-q^{j-D})=O_{q,k}(q^{-D}).
\]
Conditional on independence, their span is a uniform $k$-space and they
are a uniform ordered basis. Conditioning therefore changes the law by
$o(1)$. Apply \cref{thm:pencil} to all $m_k$ projective coefficient
points. Every limiting pattern probability in \eqref{eq:pattern} is positive,
which also proves eventual existence.
\end{proof}

\begin{remark}[Relation to matrix-space density questions]
The case of unrestricted two-dimensional matrix spaces has close antecedents
in \cite[Section 6]{gruica}, which proves the limiting probability
$g_\infty^{q+1}$ that every nonzero matrix in such a space is invertible.
That paper also studies alternating codes in a regime with growing field
size. Our statement fixes $q$ and the dimension of an alternating space,
lets its even matrix order grow, and determines every projective singularity
pattern. The underlying density questions and the classical marginal rank
law are established subjects, not new concepts introduced here.
\end{remark}
